\section{Carré magique quantique non semi-classique dans un coin APP}
\label{sec:cpe-qms-no-semiclasico}

Dans \(\mathbb C^4\), de base \(e_1,e_2,e_3,e_4\), considérons le tableau de
vecteurs
\[
\begin{array}{cccc}
e_1&e_2&e_3&e_4\\
e_2&e_1&e_4&e_3\\
(e_3+e_4)/\sqrt2&(e_3-e_4)/\sqrt2&(e_1+e_2)/\sqrt2&(e_1-e_2)/\sqrt2\\
(e_3-e_4)/\sqrt2&(e_3+e_4)/\sqrt2&(e_1-e_2)/\sqrt2&(e_1+e_2)/\sqrt2.
\end{array}
\]
Soit \(P_{ij}\) le projecteur de rang un sur le vecteur situé dans la
cellule \((i,j)\). Chaque ligne et chaque colonne est une base orthonormée, donc
\begin{equation}
 \sum_jP_{ij}=I_4,
 \qquad \sum_iP_{ij}=I_4.
 \label{eq:cpe-qms4}
\end{equation}
Pour le réaliser dans la fibre APP de dimension \(9\), on fixe
\begin{equation}
 A_{ij}:=P_{ij}\oplus\frac14I_5\in M_9(\mathbb C).
 \label{eq:cpe-qms9}
\end{equation}

\begin{theorem}[Non-semi-classicité certifiée]
\label{thm:cpe-qms-no-semiclasico}
La matrice \(A=(A_{ij})_{i,j=1}^4\) est un carré magique quantique et n'admet pas
une décomposition semi-classique
\[
 A_{ij}=\sum_{\pi:\,\pi(i)=j}Q_\pi,
 \qquad Q_\pi\succeq0,
 \qquad\sum_\pi Q_\pi=I_9.
\]
\end{theorem}

\begin{proof}
Les identités de ligne et de colonne découlent de
\eqref{eq:cpe-qms4} ; quatre copies de \(I_5/4\) ont pour somme \(I_5\). S'il existait
une décomposition semi-classique, sa compression au coin \(\mathbb C^4\)
serait une mesure conjointe des PVM définies par les lignes
\((P_{ij})_j\). Les mesures projectives conjointement mesurables commutent.
Cependant,
\[
 \left[
 |e_1\rangle\langle e_1|,
 \left|\frac{e_1+e_2}{\sqrt2}\right\rangle
 \left\langle\frac{e_1+e_2}{\sqrt2}\right|
 \right]
 =\frac12
 \begin{pmatrix}0&1\\-1&0\end{pmatrix}
 \oplus0_2\ne0.
\]
La contradiction prouve la non-semi-classicité. La compression fournit, en
même temps, un témoin analytique d'infaisabilité du problème semi-défini.
\end{proof}

Chaque ligne définit un canal UCP
\[
 \Phi_i:\ell^\infty_4\to M_9,
 \qquad \Phi_i(\delta_j)=A_{ij},
\]
et chaque colonne en définit un autre. Ainsi, les normalisations magiques se traduisent en
applications complètement positives aux domaines et codomaines explicites ; on n'utilise pas
un symbole \(J(\operatorname{id})\) sans avoir fixé ces types.

\section{Système conique de routes et universalité généalogique interne}
\label{sec:cpe-s-system-universalidad}

Pour chaque profil fini \(c\), soit \(G_c\) le groupoïde fini des routes TPK
ayant ce profil et soit \(A(c)=C^*(G_c)\). Au niveau matriciel \(s\), on prend le
cône
\[
 \mathcal C_s(c)=M_s(A(c))_+.
\]
Une route \(r:c\to d\), représentée par l'isométrie partielle \(v_r\),
induit l'application CP
\[
 T_r(a)=v_r^*av_r.
\]

\begin{theorem}[Système conique fonctoriel des routes TPK]
\label{thm:cpe-s-system-tpk}
Les assignations \(c\mapsto(\mathcal C_s(c))_{s\ge1}\) et
\(r\mapsto T_r\) respectent l'identité, la composition, l'amplification matricielle et la
troncature généalogique. La catégorie CPM engendrée par les fibres finies
est monoïdale compacte : le dual d'une fibre est son espace conjugué et les
coévaluations et évaluations sont les coupes et calottes canoniques. La limite
projective de ces systèmes conserve toutes les applications CP compatibles.
\end{theorem}

\begin{proof}
La positivité complète de \(a\mapsto v_r^*av_r\) découle immédiatement de sa forme
de Kraus. Pour des routes composables,
\(v_{r_2r_1}=v_{r_2}v_{r_1}\), donc
\(T_{r_2r_1}=T_{r_1}T_{r_2}\) avec la variance déclarée. Les inclusions et les
espérances conditionnelles des niveaux finis sont UCP et forment des carrés
commutatifs. La compactification de CPM en dimension finie est la construction
standard par espaces conjugués ; la compatibilité de troncature permet de
passer à la limite coordonnée par coordonnée.
\end{proof}

Soit \(\mathsf P_{\rm HMT}\) l'ensemble dénombrable des codes de routes
finies et soit \(\mathsf C_{\rm HMT}\) l'espace des états avec leur contexte
complet. L'évaluation
\[
 \operatorname{Ev}(\ulcorner r\urcorner,x)=T_r(x)
\]
conserve l'origine, la destination, le feuillet, l'orientation, la retenue, la frontière et le mot de
route.

\begin{corollary}[Universalité généalogiquement fidèle dans le domaine HMT]
\label{cor:cpe-universalidad-genealogica}
L'évaluateur \(\operatorname{Ev}\) est universel pour la classe déclarée des
transports TPK finis et de leurs limites compatibles : tout transport de cette
classe possède un code, la composition se réalise en concaténant les codes et la
réduction conserve la généalogie complète. Cette universalité interne ne s'
étend pas par décret à une classe extérieure de cibles non typée.
\end{corollary}

\section{Identification de la mémoire de Stinespring au registre local TPK}
\label{sec:cpe-stinespring-tpk}

Dans l'espérance nonadique, la dilatation de Stinespring conserve dans une base
\(\{e_a\}_{a=0}^8\) le chiffre écarté. Pour un état TPK \(x\), soit
\(\lambda(x,a)\) l'étiquette complète de la continuation par le chiffre
\(a\). On définit l'isométrie
\[
 W_x:e_a\longmapsto e_{\lambda(x,a)}
\]
de la mémoire minimale de Stinespring vers le registre local TPK\@. Soit
\(F_xe_{\lambda(x,a)}=e_a\) le foncteur d'oubli restreint à son image.

\Needspace{12\baselineskip}
\begin{theorem}[La mémoire TPK raffine la mémoire minimale de Stinespring]
\label{thm:cpe-stinespring-tpk}
Pour tout \(x\),
\[
 W_x^*W_x=I_9,
 \qquad F_xW_x=I_9.
\]
La famille \((W_x)_x\) entrelace l'inclusion d'un nouveau chiffre et les
troncatures. À la profondeur \(m\), le tenseur de ces isométries identifie
l'ancilla \((\mathbb C^9)^{\otimes m}\) au sous-registre TPK qui conserve
la suite des extensions ; les coordonnées supplémentaires de feuillet, de retenue,
de frontière et de route expliquent pourquoi la mémoire TPK est un raffinement et non une
seconde copie de l'ancilla.
\end{theorem}

\begin{proof}
Des continuations de chiffres distincts possèdent des étiquettes orthogonales, de sorte
que \(W_x\) est isométrique. La définition de \(F_x\) prouve la seconde
identité. La composition des extensions concatène les étiquettes dans le même
ordre que le produit tensoriel des ancillas ; d'où la naturalité sous
troncature et l'affirmation à profondeur arbitraire.
\end{proof}
